% 乔治·基斯佳科夫斯基（George Kistiakowsky）（综述）
% license CCBYSA3
% type Wiki

本文根据 CC-BY-SA 协议转载翻译自维基百科\href{https://en.wikipedia.org/wiki/George_Kistiakowsky}{相关文章}。

\begin{figure}[ht]
\centering
\includegraphics[width=6cm]{./figures/a36c1f246ada9580.png}
\caption{} \label{fig_QZJS_1}
\end{figure}
乔治·博格达诺维奇·基斯佳科夫斯基（俄语：Георгий Богданович Кистяковский；乌克兰语：Георгій Богданович Кістяківський，罗马化：Heorhii Bohdanovych Kistiakivskyi；1900年12月1日［旧历11月18日］—1982年12月7日）是一位乌克兰裔美国物理化学家，曾任哈佛大学教授，参与“曼哈顿计划”，并在后期担任美国总统德怀特·D·艾森豪威尔的科学顾问。

他出生于当时俄帝国境内的博亚尔卡（Boyarka）\(^\text{[1]}\)，出身于“一个古老的乌克兰哥萨克家族，该家族在俄国革命前是知识精英阶层的一部分”\(^\text{[2]}\)。在俄国内战期间，基斯佳科夫斯基逃离祖国，前往德国，在柏林大学师从马克斯·博登斯坦（Max Bodenstein），获得物理化学博士学位。1926年他移居美国，1930年加入哈佛大学任教，并于1933年成为美国公民。

第二次世界大战期间，基斯佳科夫斯基担任国家国防研究委员会（NDRC）爆炸物开发部门负责人，同时兼任爆炸物研究实验室（ERL）技术主任，负责监督包括RDX与HMX在内的新型炸药的研制。他还参与了爆炸的流体动力学理论研究以及成形装药技术的发展。1943年10月，他被聘为“曼哈顿计划”的顾问，不久后接管X部门，负责研发用于内爆型核武器的爆炸透镜。1945年7月，他亲眼见证了“三位一体”核试验中的首次原子爆炸。数周后，另一枚内爆型核武器“胖子”（Fat Man）被投放至长崎。

1962年至1965年间，基斯佳科夫斯基担任美国国家科学院科学、工程与公共政策委员会（COSEPUP）主席，并于1965年至1973年间出任该院副院长。由于反对越南战争，他切断了与政府的所有联系，并积极投身反战组织“宜居世界委员会”（Council for a Livable World），并于1977年出任该组织主席。
